\section*{Chapitre \ref{ch:b3:homogeneous-catalysis} --- La catalyse homogène dans l'industrie}

\begin{solution}{exo:b3:homogeneous-catalysis:1}
$\ce{RhCl(PPh3)3}$~: 16, Rh(I) $d^8$. $\ce{RhCl(PPh3)2}$~: 14, Rh(I). $\ce{RhClH2(PPh3)2}$~: 16, Rh(III) $d^6$.
Complexe de l'alcène~: 18, Rh(III). Hydrure d'alkyle~: 16, Rh(III).
L'élimination réductrice ramène à 14, Rh(I).
\end{solution}

\begin{solution}{exo:b3:homogeneous-catalysis:2}
Le butanal, $\ce{CH3CH2CH2CHO}$ (linéaire), et le 2-méthylpropanal, $\ce{(CH3)2CHCHO}$ (ramifié).
\end{solution}

\begin{solution}{exo:b3:homogeneous-catalysis:3}
Le cyclopent-3-ène-1,1-dicarboxylate de diéthyle, un cycle à cinq chaînons,
avec perte d'éthène~: les deux groupes $\ce{CH2}$ terminaux partent sous forme de $\ce{C2H4}$ et les deux carbones CH
internes se relient.
\end{solution}

\begin{solution}{exo:b3:homogeneous-catalysis:4}
Le (\textit{E})-3-phénylprop-2-énoate de méthyle (cinnamate de méthyle)~:
l'aryle s'additionne sur le carbone terminal, et l'élimination d'hydrure en $\beta$, syn, à partir du conformère le plus stable,
donne l'alcène \textit{E}.
\end{solution}

\begin{solution}{exo:b3:homogeneous-catalysis:5}
L'addition oxydante est lente et $\ce{[Rh(CO)2I2]-}$ est l'état de repos, donc
$v = k[\mathrm{Rh}]_{\mathrm{tot}}[\ce{CH3I}]$. Le méthanol ne fait que régénérer l'iodure de méthyle par une réaction rapide avec HI~;
l'iodure total introduit fixe $[\ce{CH3I}]$, si bien que la vitesse ne dépend pas de la
concentration en méthanol tant qu'il reste assez de méthanol pour
reconvertir HI en iodure de méthyle.
\end{solution}

\begin{solution}{exo:b3:homogeneous-catalysis:6}
L'acide phénylboronique avec le 4-bromotoluène, ou l'acide
4-méthylphénylboronique avec le bromobenzène. Les deux fonctionnent~; la
première voie utilise l'acide boronique le moins cher et un bromure d'aryle
stable et courant.
\end{solution}

\begin{solution}{exo:b3:homogeneous-catalysis:7}
Symétrie $C_2$~: isotactique (les deux sites sélectionnent la même face). Symétrie $C_s$~:
syndiotactique (les sites images l'un de l'autre dans un miroir alternent). $\ce{Cp2ZrCl2}$ non ponté~: atactique (sites
achiraux).
\end{solution}

\begin{solution}{exo:b3:homogeneous-catalysis:8}
$\mathrm{TON} = 0.98/0.0010 = 980$~; $\mathrm{TOF} = 980/\qty{2.0}{h} = \qty{490}{h^{-1}}$ en moyenne.
\end{solution}

\begin{solution}{exo:b3:homogeneous-catalysis:9}
Le motif de répétition est $\ce{-[CH=CH-C5H8]-}$, les deux groupes CH étant portés par les carbones 1 et 3 d'un
cycle cyclopentane (poly(1,3-cyclopentylènevinylène)). La libération de la
tension du cycle bicyclique rend $\Delta_rH^\circ$ fortement négatif, ce qui l'emporte sur la
perte d'entropie de translation~; aucun gaz ne se forme, si bien que la
force motrice est enthalpique.
\end{solution}

\begin{solution}{exo:b3:homogeneous-catalysis:10}
$\bar n_n = 1/(1 - p) = 2000$~; $M_w/M_n = 1 + p = 1.9995 \approx 2.0$. Un \omterm{def:b3:homogeneous-catalysis:ziegler-natta}{catalyseur de Ziegler--Natta} comporte
plusieurs sortes de sites, chacun avec son propre $p$~; la somme de plusieurs distributions de Schulz--Flory
de moyennes différentes est plus large (dispersité souvent de 4 à 8).
\end{solution}

\begin{solution}{exo:b3:homogeneous-catalysis:11}
À basse pression, $K_1K_2[\ce{H2}] \ll [\mathrm L] + K_1$~: $v \approx kK_1K_2[\mathrm{Rh}]_{\mathrm{tot}}[\text{alcène}][\ce{H2}]/([\mathrm L] + K_1)$,
du premier ordre par rapport à $\ce{H2}$. À haute pression, le terme $K_1K_2[\ce{H2}]$ domine~: $v \to k[\mathrm{Rh}]_{\mathrm{tot}}[\text{alcène}]$,
indépendante de $\ce{H2}$. $[\mathrm L]$ n'apparaît qu'au dénominateur~: l'ajout de phosphine diminue $v$
en ramenant le rhodium vers le précatalyseur saturé.
\end{solution}

\begin{solution}{exo:b3:homogeneous-catalysis:12}
L'insertion qui place le rhodium sur le carbone interne construit un alkyle
secondaire à côté d'un métal encombré~; les phosphines encombrées (et leur
excès, qui en maintient deux sur le métal) rendent cet état de transition
plus défavorable que celui qui donne l'alkyle primaire, lequel conduit à
l'aldéhyde linéaire. C'est le butanal linéaire qui est converti, par
crotonisation et hydrogénation, en l'alcool ramifié en $C_8$ utilisé dans les
plastifiants~; le 2-méthylpropanal est un sous-produit de moindre valeur.
\end{solution}

\begin{solution}{pb:b3:homogeneous-catalysis:1}
\textbf{1.} Ir(I), $d^8$~: $9 + 4 + 2 + 1 = 16$ électrons (décompte neutre, en tenant compte de la charge).

\textbf{2.} $\ce{[CH3Ir(CO)2I3]-}$, 18 électrons, Ir(III).

\textbf{3.} $\ce{[CH3Ir(CO)3I2]}$, 18 électrons, Ir(III).

\textbf{4.} $\ce{[CH3COIr(CO)2I2]}$, 16 électrons, Ir(III).

\textbf{5.} $\ce{[CH3COIr(CO)2I3]-}$ (18) élimine $\ce{CH3COI}$, ce qui régénère $\ce{[Ir(CO)2I2]-}$.

\textbf{6.} $\ce{CH3COI + H2O -> CH3COOH + HI}$~; $\ce{CH3OH + HI -> CH3I + H2O}$.

\textbf{7.} $\ce{CH3OH + CO -> CH3COOH}$.

\textbf{8.} Avec l'iridium, c'est l'insertion migratoire qui est lente~;
l'addition oxydante de l'iodure de méthyle, cinétiquement déterminante avec
le rhodium, est bien plus rapide sur l'iridium, plus riche en électrons.

\textbf{9.} Ils arrachent un iodure à $\ce{[CH3Ir(CO)2I3]-}$, ce qui donne le tricarbonyle neutre, dans lequel
le métal rétrocède moins et le méthyle migre plus vite sur un CO plus
électrophile.

\textbf{10.} Une dépendance inverse vis-à-vis de la concentration en iodure
libre~: l'iodure ramène le pré-équilibre vers le complexe anionique lent.

\textbf{11.} Le méthanol n'intervient qu'après l'étape cinétiquement
déterminante, converti en iodure de méthyle par une réaction rapide avec
HI.

\textbf{12.} $M(\ce{CH3COOH}) = \qty{60.1}{g/mol}$~: $n = \qty{5.0e11}{g}/\qty{60.1}{g/mol} = \qty{8.3e9}{mol}$ par an.

\textbf{13.} $\qty{8.32e9}{mol}/\qty{8000}{h} = \qty{1.04e6}{mol/h}$.

\textbf{14.} $\qty{2.0e5}{L} \times \qty{5.0e-3}{mol/L} = \qty{1.0e3}{mol}$ d'iridium, $\qty{1.0e3}{mol} \times \qty{192.2}{g/mol} \approx \qty{190}{kg}$.

\textbf{15.} $\mathrm{TOF} = \qty{1.04e6}{mol/h}/\qty{1.0e3}{mol} = \qty{1.0e3}{h^{-1}}$, soit environ \qty{0.29}{s^{-1}}.

\textbf{16.} $\mathrm{TON} = \qty{8.3e9}{mol}/\qty{1.0e3}{mol} = \num{8.3e6}$ par an.

\textbf{17.} $\qty{5.0e8}{kg}/(\qty{8000}{h} \times \qty{200}{m^3}) \approx \qty{310}{kg.m^{-3}.h^{-1}}$.

\textbf{18.} $\ce{CO + H2O -> CO2 + H2}$.

\textbf{19.} $1/0.94 = \qty{1.064}{mol}$ de CO par mole d'acide acétique.

\textbf{20.} $1.064 - 1 = \qty{0.064}{mol}$ de \ce{CO2} par mole d'acide acétique.

\textbf{21.} Une tonne représente $\qty{1.0e6}{g}/\qty{60.1}{g/mol} = \qty{1.66e4}{mol}$~; $\qty{1.66e4}{mol} \times 0.0638 \times \qty{44.0}{g/mol} \approx \qty{47}{kg}$ de \ce{CO2}.

\textbf{22.} L'eau est un réactif de la réaction du gaz à l'eau~: moins
d'eau la ralentit et économise l'énergie de séchage de l'acide. L'eau reste
nécessaire pour hydrolyser l'iodure d'acétyle et pour maintenir le
catalyseur actif et dissous~; avec le rhodium, une faible teneur en eau
fait précipiter l'iodure de rhodium.

\textbf{23.} Le dihydrogène~: il s'accumule dans le gaz, abaisse la
pression partielle de CO (il faut donc purger du gaz, ce qui fait perdre du
CO), et peut hydrogéner des intermédiaires en sous-produits comme l'acide
propanoïque.

\textbf{24.} La fréquence de cycles du catalyseur à l'iridium vaut environ $\qty{1.0e3}{h^{-1}}$~: chaque atome d'iridium produit
un millier de molécules d'acide acétique par heure, quelque huit millions
par an.
\end{solution}
